\question[3]
Erläutere, warum jedes Gerät in einem Netzwerk eine eindeutige Adresse benötigt.
 Nenne dazu zwei Beispiele für Adressierungssysteme aus dem Alltag 
 außerhalb von Computernetzen und erkläre kurz die Parallele zur IP-Adresse.

\begin{solutionbox}{6cm}

Eine eindeutige Adresse ist nötig, damit Datenpakete an genau
den richtigen Empfänger zugestellt werden können. Ohne
eindeutige Adressierung wüsste kein Gerät, für wen eine
Nachricht bestimmt ist, oder mehrere Geräte könnten die
gleiche Nachricht fälschlich erhalten.

Beispiele: \\
- Handynummer: eindeutig einer SIM-Karte zugeordnet,
  vergleichbar mit der IP-Adresse eines Geräts. \\
- E-Mail-Adresse: eindeutig einem Postfach zugeordnet,
  ermöglicht gezielte Zustellung von Nachrichten, so wie die
  IP-Adresse die gezielte Zustellung von Datenpaketen
  ermöglicht.

\end{solutionbox}